% META: {"id": "TCOMPL-AIR-EX-029", "chapitre": "TCOMPL-CALCULS-AIRES", "type_objet": "exercice", "capacites": ["TCOMPL-AIR-C5"], "capacites_codes": ["C5"], "methodes": ["M4"], "parcours": 1, "competences": ["communiquer","raisonner"], "duree_min": 12, "mode_creation": "ex_nihilo", "sources_inspiration": [], "fichier_tex": "chapitres/TCOMPL-CALCULS-AIRES/exercices/TCOMPL-AIR-EX-029.tex", "corrige_tex": "chapitres/TCOMPL-CALCULS-AIRES/corriges/TCOMPL-AIR-CO-029.tex", "status": "approved"}
% BEGIN-VERIFY
% from sympy import *
% x = symbols('x')
% f = x**4 - 4*x**3 + 6*x**2
% fpp = diff(f, x, 2)
% assert expand(fpp) == 12*x**2 - 24*x + 12
% assert factor(fpp) == 12*(x-1)**2
% # f'' >= 0 partout, pas de changement de signe => pas d'inflexion
% assert fpp.subs(x, 1) == 0
% assert fpp.subs(x, 0) == 12
% assert fpp.subs(x, 2) == 12
% END-VERIFY
\begin{exercice}{TCOMPL-AIR-EX-029}{1}{12}
Soit $f(x) = x^4 - 4x^3 + 6x^2$. La courbe admet-elle des points d'inflexion ?

\end{exercice}
